\subsection{The special return and its old bound}\label{sec:special-details}
For a special step let $R$ be the actual retained old occurrences, and
let $\pi$ be the retention correspondence into the completed target.
Define
\begin{equation}\label{eq:7-5}
 D=\{u\in R:\Gamma_M(u)\ne\Gamma_{M[n]}(\pi(u))\}.
\end{equation}
Both consumers are computed on their own original completed tables;
the comparison uses actual occurrences under retention, not new program
identifiers. If $D\ne\varnothing$, let $u_*$ be its last member in
the original depth-first item order and put $b=\ell_M(u_*)$.
The permission for this special operation is bounded by this old $b$.
The construction above obtains every actual upward field strictly below
$b$, which is stronger than the allowed weak bound.

To see why this bound is available, traverse the unique changed
containing chain from the inside outward. Before reaching the first
member of $D$, candidate retention or old persistent coverage keeps
each actual incoming consumer. Complete-tail deletion and recursive
prefix supply there use the old whole-entry downward statement: the
outer label and Gen can remain, its old graph remains exact, and D4
restricts to the new last direct column. Any newly required low
support is read from the retained complete packages and certified by
their birth records. No upward operation has yet been used.

The first changed retained entry encountered is $u_*$. All genuinely
smaller internal targets at that call are now closed. Thus the old
entry gives a logical high Point using its retained D4 work model.
Apply the common-source existence statement to actual sources before
$u_*$, choosing the entire new Point below its old label. This uses
the closed high configuration as an existence witness; it does not
first raise a fresh object inside the old private height and then
raise the enclosing package. Any inner auxiliary values used to
exhibit that configuration are quantified anew with the whole Point.

Continue the outside returns with the shared source bounds, preserving
all old complete-prefix readers. If the special rule makes several
same-parent copies, finish them in their actual order; a future copy
whose private height is still open is not a closed proper right carrier.
The larger outside old bounds are unnecessary: every new upward field
still lies below $b$. If $D$ is empty, this changed-entry step is
absent and the entire special construction has no upward operation.
Finally replay and minimal support identify the target with its
original completion. The deepest old tail is deleted and no field
reaches the old maximum, proving strict maximum descent.

This separate argument uses neither the general connection stages nor
the adjacent-index permission criterion. A change of a retained old
height by itself is not a witness for (\ref{eq:7-5}); an actual change of its
complete consumer is required.
